% !TeX root = ../../金融经济学_王江_讲义及习题.tex
% 来源：教材印刷第 108—109 页（PDF 第 123—124 页），练习 7.1—7.5。

\begin{problem}
\WJSevenExercise{ex:wj-seven-1}{} 令 $R(w)\equiv -wu''(w)/u'(w)$ 而
\[
    a(w)=\operatorname*{arg\,max}_{a} E\bigl[u(w+a\tilde{r})\bigr]
\]
为相对于效用函数 $u(\cdot)$ 的最优投资选择。假设 $w>0$、$u'(\cdot)>0$、$u''(\cdot)\le 0$ 且 $E[\tilde{r}]>0$。证明以下命题：
\begin{enumerate}[label=(\alph*)]
    \item 如果 $\lambda>1$，给定效用函数 $u(\cdot)$ 且 $R'(w)\ge 0$。定义
    \[
        u_{1}(w)\equiv u(\lambda w),\qquad u_{2}(w)\equiv u(w).
    \]
    那么，$R_{1}(w)\ge R_{2}(w)$。
    \item 给定以上结论，我们有 $a_{1}(w)\le a_{2}(w)$，其中 $a_{1}(w)$ 和 $a_{2}(w)$ 分别是相对于 $u_{1}$ 和 $u_{2}$ 的最优投资选择。
    \item 定义 $\widehat{a}(w)\equiv a(w)/w$。那么，对于 $w_{1}\le w_{2}$ 有 $\widehat{a}(w_{1})\ge\widehat{a}(w_{2})$。（提示：令 $\lambda=w_{2}/w_{1}$。）
\end{enumerate}
最后一个结论说的是，对于具有单调递增的相对风险厌恶的投资者，他在风险资产上的投资额占他财富的比例随着财富的增加而减少。
\end{problem}
\begin{solution}[label-content={证明},blank-type=manual,blank-vsep=11\baselineskip]
记初始财富为 $w$ 时的可行投资区间为 $D(w)$。它包含不投资的选择 $a=0$，并要求可行头寸在各可能状态下都使期末财富为正，以便在正财富区间上比较风险厌恶。对 $i=1,2$，用同一可行集求解
\[
 \begin{aligned}
  &\max_a\quad \mathrm E[u_i(w+a\tilde r)]\\
  &\text{s.t.}\quad a\in D(w).
 \end{aligned}
\]

假设期望有限、最优投资额存在且唯一，并且在区间内部可将导数移入期望。最优额是否存在取决于收益分布、效用和可行集。例如，若效用线性、收益确定且为正，投资额又无上限，期望效用就会随投资额一直上升。
\begin{enumerate}[label=(\alph*)]
\item 由链式法则，
\[
 u_1'(z)=\lambda u'(\lambda z),\qquad
 u_1''(z)=\lambda^2u''(\lambda z).
\]

因此
\[
 R_1(z)=-\frac{z u_1''(z)}{u_1'(z)}
 =-\frac{\lambda z u''(\lambda z)}{u'(\lambda z)}
 =R(\lambda z)\ge R(z)=R_2(z),
\]
其中用到 $z>0$、$\lambda>1$ 和 $R'\ge0$。

\item 要比较两个最优投资额，先比较同一投资额下增加一单位投资带来的边际效用。记绝对风险厌恶系数为 $A_i(z)=-u_i''(z)/u_i'(z)$。由 (a)，
\[
 A_1(z)=\frac{R_1(z)}{z}\ge\frac{R_2(z)}{z}=A_2(z).
\]

令 $h(z)=u_1'(z)/u_2'(z)>0$，则
\[
 \frac{h'(z)}{h(z)}=-A_1(z)+A_2(z)\le0,
\]

所以 $h$ 非增。对给定的初始财富 $w$，记
\[
 F_i(a)=E[u_i(w+a\tilde r)],\qquad
 F_i'(a)=E[u_i'(w+a\tilde r)\tilde r].
\]

当可行的 $a<0$ 时，由凹效用的切线界，$F_i(a)\le u_i(w)+u_i'(w)aE[\tilde r]<F_i(0)$，所以两个最优额都非负。取任意正的内点投资额 $a$。当 $\tilde r\ge0$ 时，$h(w+a\tilde r)\le h(w)$；当 $\tilde r<0$ 时，不等式反向，但乘以 $\tilde r$ 后又反向。因此逐状态有
\[
 h(w+a\tilde r)\tilde r\le h(w)\tilde r.
\]

乘以正数 $u_2'(w+a\tilde r)$ 并取期望，得到
\[
 \begin{aligned}
 F_1'(a)
 &=E[h(w+a\tilde r)u_2'(w+a\tilde r)\tilde r]\\
 &\le h(w)F_2'(a).
 \end{aligned}
\]

若 $a_1(w)>a_2(w)$，在两者之间取 $a$。由凹性和最优点的唯一性，$F_1'(a)>0$，因为右侧的 $a_1(w)$ 给出更高的 $F_1$；同理，左侧的 $a_2(w)$ 给出更高的 $F_2$，所以 $F_2'(a)<0$。这与上式矛盾，故 $a_1(w)\le a_2(w)$。取导数的位置始终在两个最优额之间，因此最优额位于可行区间端点时也适用。

\item 若 $w_1=w_2$，结论显然。若 $0<w_1<w_2$，令 $\lambda=w_2/w_1>1$。在初始财富 $w_1$ 下，
\[
 E[u_1(w_1+a\tilde r)]
 =E[u(w_2+\lambda a\tilde r)].
\]

令 $b=\lambda a$，则 $w_2+b\tilde r=\lambda\cdot(w_1+a\tilde r)$。若可行性只要求期末财富为正，两边同时为正，故 $a\in D(w_1)$ 当且仅当 $b\in D(w_2)$。若另有头寸限制，则这里要求这些限制也随财富同比例缩放，即 $D(w_2)=\lambda D(w_1)$。在这个条件下，$b$ 恰好遍历财富为 $w_2$ 时的全部可行投资额。因此
\[
 a_1(w_1)=\frac{a(w_2)}{\lambda},\qquad a_2(w_1)=a(w_1).
\]

由 (b) 得
\[
 \frac{a(w_2)}{\lambda}\le a(w_1)
 \quad\Longrightarrow\quad
 \boxed{\frac{a(w_2)}{w_2}\le\frac{a(w_1)}{w_1}}.
\]

因此相对风险厌恶随财富非减时，最优风险投资比例随财富非增。
\end{enumerate}
\end{solution}

\begin{problem}
\WJSevenExercise{ex:wj-seven-2}{} 参与者的初始财富为 $w$ 且他的绝对风险厌恶系数为常数 $a$。现在他必须承担风险 $\tilde{g}$，$\tilde{g}$ 是一公平赌博（即，他的财富变成 $w+\tilde{g}$）。当 $\tilde{g}$ 具有如下分布时计算确定性等价：
\begin{enumerate}[label=(\alph*)]
    \item 取值为 $-b$ 和 $b$ 的二项分布；
    \item 区间 $[-c,c]$ 上的均匀分布；
    \item 均值为 0、标准差为 $\sigma$ 的正态分布。
    \item 讨论在以上情形下，确定性等价如何依赖于初始财富 $w$。从中你能得到什么结论？并解释你的结论。
\end{enumerate}
\end{problem}
\begin{solution}[blank-type=manual,blank-vsep=9\baselineskip]
常绝对风险厌恶 $a>0$ 时，可取 $u(x)=-e^{-ax}$，其绝对风险厌恶系数 $-u''(x)/u'(x)$ 恒为 $a$。这个效用在整个实数轴上有定义，因此也适用于期末财富可能为负的正态情形。按风险溢价的定义\ifdefined\WJHandoutKnowledge（\WJRef{def:wj-seven-4}{定义~}[（风险溢价）]）\fi，
\[
 \mathrm E[u(w+\tilde g)]=u(w-\pi)=u(w+C),\qquad
 C=-\pi,\qquad CE=w-\pi=w+C.
\]

其中 $\pi$ 是为消除风险愿意支付的金额，$C$ 是赌博净额的确定性等价，$CE$ 是期末财富的确定性等价。代入负指数效用，
\[
 -e^{-a\cdot(w+C)}=-e^{-aw}\mathrm E[e^{-a\tilde g}]
\]
约去共同因子 $-e^{-aw}$ 后取对数，得到
\[
 C=-\frac1a\ln\mathrm E[e^{-a\tilde g}].
\]

若 $a=0$，可取线性效用，三种公平赌博都有 $C=\mathrm E[\tilde g]=0$。以下计算 $a>0$ 的情形。
\begin{enumerate}[label=(\alph*)]
\item 公平性使 $-b,b$ 各以 $1/2$ 的概率出现，因此
\[
 \mathrm E[e^{-a\tilde g}]=\tfrac12(e^{ab}+e^{-ab})=\cosh(ab),
 \qquad C=-\frac1a\ln\cosh(ab).
\]
\item 均匀分布下，
\[
 \mathrm E[e^{-a\tilde g}]=\frac1{2c}\int_{-c}^ce^{-ax}\,\mathrm dx
 =\frac{\sinh(ac)}{ac},\qquad
 C=-\frac1a\ln\frac{\sinh(ac)}{ac}.
\]
\item 正态密度为 $\frac1{\sigma\sqrt{2\pi}}e^{-x^2/(2\sigma^2)}$。配方得
\[
 -ax-\frac{x^2}{2\sigma^2}
 =-\frac{(x+a\sigma^2)^2}{2\sigma^2}+\frac{a^2\sigma^2}{2},
\]

配方后的正态密度以 $-a\sigma^2$ 为均值、方差仍为 $\sigma^2$，积分为 $1$。因此 $\mathrm E[e^{-a\tilde g}]=e^{a^2\sigma^2/2}$，从而
\[
 C=-\frac{a\sigma^2}{2},\qquad\pi=\frac{a\sigma^2}{2}.
\]
\item 固定赌博的金额及其分布，初始财富因子 $-e^{-aw}$ 已在求 $C$ 时约去。因此净额 $C$ 和消除风险的支付意愿 $\pi=-C$ 均与 $w$ 无关。期末财富的确定性等价 $CE=w+C$ 则随 $w$ 一比一增加。
\end{enumerate}
\end{solution}

\begin{problem}
\WJSevenExercise{ex:wj-seven-3}{} 参与者的初始财富为 $w$ 且他的绝对风险厌恶系数为 1。现在他必须承担与他的财富成比例的风险 $w\tilde{g}$，$w\tilde{g}$ 是一公平赌博。当 $\tilde{g}$ 具有如下分布时，计算确定性等价占初始财富的比例——即相对确定性等价：
\begin{enumerate}[label=(\alph*)]
    \item 取值为 $-b$ 和 $b$ 的二项分布，其中 $0<b<1$；
    \item 区间 $[-c,c]$ 上的均匀分布，其中 $0<c<1$；
    \item 讨论在以上情形下，相对确定性等价如何依赖于初始财富 $w$。从中你能得到什么结论？并解释你的结论。
\end{enumerate}

% 原书印刷第 109 页；参考答案采用对数效用，与原题的“绝对”条件不一致，保留原题并分别处理。
\end{problem}
\begin{solution}[blank-type=manual,blank-vsep=9\baselineskip]
绝对风险厌恶系数恒为 $1$，可取 $u(x)=-e^{-x}$。沿用上一题的净额口径，相对确定性等价为 $C_R=C/w$，其中 $w>0$。将 $a=1$、净风险 $w\tilde g$ 代入上一题的公式，得到 $C=-\ln\mathrm E[e^{-w\tilde g}]$。若按期末财富总额计算，对应的比例是 $CE/w=1+C_R$。
\begin{enumerate}[label=(\alph*)]
\item 对公平二项风险，
\[
 C=-\ln\cosh(wb),\qquad
 C_R=-\frac{\ln\cosh(wb)}{w}.
\]
\item 对均匀风险，
\[
 C=-\ln\frac{\sinh(wc)}{wc},\qquad
 C_R=-\frac1w\ln\frac{\sinh(wc)}{wc}.
\]
\item 取 $0<w_1<w_2$，记 $\lambda=w_1/w_2$。两种分布都不是确定支付，因此 $e^{-w_2\tilde g}$ 也不是常数。由于 $x^\lambda$ 在正实数上严格凹，严格 Jensen 不等式给出
\[
 \mathrm E[e^{-w_1\tilde g}]
 =\mathrm E\!\left[(e^{-w_2\tilde g})^\lambda\right]
 <\bigl(\mathrm E[e^{-w_2\tilde g}]\bigr)^\lambda.
\]

两边取对数再乘以 $-1/w_1$，不等号反向。由 $\lambda/w_1=1/w_2$，得到
\[
 C_R(w_1)=-\frac{\ln\mathrm E[e^{-w_1\tilde g}]}{w_1}
 >-\frac{\ln\mathrm E[e^{-w_2\tilde g}]}{w_2}=C_R(w_2).
\]

所以财富越高，$C_R$ 越低，愿意为消除风险而支付的财富比例 $-C_R$ 越高。这里赌博的金额 $w\tilde g$ 会随财富放大，比较的是同一比例的风险。这与 $R(w)=wA(w)=w$ 随财富上升一致。
\end{enumerate}

若改为相对风险厌恶恒为 $1$，可取对数效用。此时 $\ln(w+C)=\mathrm E[\ln(w\cdot(1+\tilde g))]=\ln w+\mathrm E[\ln(1+\tilde g)]$。解得 $C_R=\exp(\mathrm E[\ln(1+\tilde g)])-1$，代入两种分布有
\[
 C_R=
 \begin{cases}
  \sqrt{1-b^2}-1,&\text{公平二项风险},\\[3pt]
  \displaystyle\exp\!\left(
  \frac{(1+c)\ln(1+c)-(1-c)\ln(1-c)-2c}{2c}\right)-1,
  &\text{均匀风险}.
 \end{cases}
\]

在这一对照情形下，$C_R$ 与 $w$ 无关，期末财富的确定性等值 $CE=w\cdot(1+C_R)$ 与 $w$ 成比例。
\end{solution}

\begin{problem}
\WJSevenExercise{ex:wj-seven-4}{} 参与者的初始财富为 $w$ 且相对风险厌恶系数为 2。现在他要承担风险 $\tilde{g}$，$\tilde{g}$ 是有两个可能取值 $-b$ 和 $b$ 的公平赌博，$0<b<w$。我们能用两种方式定义风险溢价：
\begin{itemize}
    \item $E[u(w+\tilde{g})]=u(w-\pi)$；
    \item $E[u(w+\tilde{g}+\pi)]=u(w)$。
\end{itemize}
用这两种定义计算风险溢价。它们相等吗？解释所得到的结论。
\end{problem}
\begin{solution}[blank-type=manual,blank-vsep=8\baselineskip]
相对风险厌恶系数恒为 $2$，所以
\[
 -\frac{x u''(x)}{u'(x)}=2
 \quad\Longrightarrow\quad
 \frac{\mathrm d}{\mathrm dx}\ln u'(x)=-\frac2x.
\]

积分得 $u'(x)=Kx^{-2}$（$K>0$），故可取
\[
 u(x)=-\frac1x,\qquad x>0.
\]

公平二项赌博的两个结果各以 $1/2$ 的概率发生。用 $\pi_1$ 表示为了消除风险愿意支付的金额。由 $\mathrm E[u(w+\tilde g)]=u(w-\pi_1)$，
\[
 -\frac12\left(\frac1{w+b}+\frac1{w-b}\right)
 =-\frac{w}{w^2-b^2}=-\frac1{w-\pi_1}.
\]

所以 $w-\pi_1=(w^2-b^2)/w>0$，从而
\[
  \pi_1=\frac{b^2}{w}.
\]

再用 $\pi_2$ 表示接受风险所需的补偿。由 $\mathrm E[u(w+\tilde g+\pi_2)]=u(w)$，
\[
 -\frac12\left(\frac1{w+\pi_2+b}+\frac1{w+\pi_2-b}\right)
 =-\frac{w+\pi_2}{(w+\pi_2)^2-b^2}=-\frac1w.
\]

令 $z=w+\pi_2$，得到 $z^2-wz-b^2=0$。取满足正财富约束的正根，得
\[
 z=\frac{w+\sqrt{w^2+4b^2}}2>w>b,
 \qquad
 \pi_2=\frac{\sqrt{w^2+4b^2}-w}{2}.
\]

把第一种风险溢价写成初始财富的函数 $\pi_1(z)$，赌博幅度 $b$ 固定。按第一种定义，在补偿后的财富 $z=w+\pi_2$ 下承担同一风险，其期望效用等于 $u(z-\pi_1(z))$。按第二种定义，这个期望效用也等于 $u(w)$。因为 $u$ 严格递增，$z-\pi_1(z)=w$，所以
\[
 \pi_2=\pi_1(w+\pi_2)=\frac{b^2}{w+\pi_2}<\frac{b^2}{w}=\pi_1(w).
\]

在本题中，绝对风险厌恶 $A(x)=2/x$ 随财富增加而下降。补偿先提高了财富，因此在较高的财富水平上衡量同一笔风险，溢价更小。把 $\pi_2$ 的根式有理化，可写成
\[
 \pi_2=\frac{b^2}{w}\cdot\frac{2}{1+\sqrt{1+4(b/w)^2}}.
\]

固定 $w$，当 $b/w$ 趋于零时，第二个因子趋于 $1$，所以两种溢价关于赌博幅度 $b$ 的二阶主导项都是 $b^2/w$。

\end{solution}

\Needspace{26\baselineskip}
\begin{problem}
\WJSevenExercise{ex:wj-seven-5}{} 参与者的初始财富为 $w$ 且是严格风险厌恶的。$\tilde{g}_{1}$ 和 $\tilde{g}_{2}$ 是两个独立的公平赌博，具有相同的、定义于 $\{-b,b\}$ 上的二项分布，其中 $0<b<w/2$。
\begin{enumerate}[label=(\alph*)]
    \item 假设他必须承担风险 $\tilde{g}_{1}$，即他的财富变成 $w+\tilde{g}_{1}$。在这种情况下，他的期望效用记为 $V_{1}=E[u(w+\tilde{g}_{1})]$。证明风险使得他的情况恶化，即 $V_{1}<u(w)$。
    \item 现在假设他必须承担的是分散化的风险 $\tilde{g}=\frac12(\tilde{g}_{1}+\tilde{g}_{2})$。记这种情形下的期望效用为 $V_{2}=E[u(w+\tilde{g})]$。$V_{2}$ 是否总是高于 $V_{1}$？
\end{enumerate}
\end{problem}
\begin{solution}[blank-type=manual,blank-vsep=7\baselineskip]
\begin{enumerate}[label=(\alph*)]
\item \textbf{证明}\quad 由公平性，$-b$ 和 $b$ 各以 $1/2$ 的概率发生。严格风险厌恶意味着伯努利效用 $u$ 严格凹\ifdefined\WJHandoutKnowledge（见\WJRef{thm:wj-seven-3}{定理~}[（风险厌恶与效用凹性）]）\fi。由于 $b>0$，两个期末财富 $w+b$ 和 $w-b$ 均为正且不同，严格 Jensen 不等式给出
\[
 V_1=\tfrac12u(w+b)+\tfrac12u(w-b)
 <u\left(\tfrac12(w+b)+\tfrac12(w-b)\right)=u(w).
\]

\item \textbf{解答}\quad 两个赌博独立，四种输赢组合各有 $1/4$ 的概率。两次都赢或都输时，平均风险分别为 $b$ 或 $-b$。一赢一输的两条路径上，盈亏相抵，合计有 $1/2$ 的概率得到零。因此
\[
 \Pr(\tilde g=b)=\tfrac14,\quad
 \Pr(\tilde g=0)=\tfrac12,\quad
 \Pr(\tilde g=-b)=\tfrac14.
\]

于是
\[
 \begin{aligned}
 V_2&=\tfrac14u(w+b)+\tfrac12u(w)+\tfrac14u(w-b)\\
 &=\tfrac12V_1+\tfrac12u(w),
 \end{aligned}
\]

从而
\[
  V_2-V_1=\tfrac12\bigl[u(w)-V_1\bigr]>0.
\]

所以在题设条件下，总有 $V_2>V_1$。分散化改善了处境，但平均后的风险仍非退化，故仍有 $V_2<u(w)$；完整排序为 $V_1<V_2<u(w)$。
\end{enumerate}
\end{solution}

